\documentclass[11pt,a4paper]{article}

\usepackage[T1]{fontenc}
\usepackage{lmodern}
\usepackage{microtype}
\usepackage[a4paper,margin=26mm]{geometry}
\usepackage{amsmath,amssymb,amsthm,mathtools}
\usepackage[hidelinks]{hyperref}
\hypersetup{
  pdftitle={The Short-Time Remainder Is One Order Higher},
  pdfsubject={Vanishing of the order 2m+2 term in short-time hypocoercivity}
}

\newtheorem{proposition}{Proposition}
\newtheorem{lemma}{Lemma}
\DeclareMathOperator{\Log}{Log}
\setlength{\parindent}{0pt}
\setlength{\parskip}{0.6em}
\allowdisplaybreaks
\emergencystretch=2em

\title{\bfseries The Short-Time Remainder Is One Order Higher}
\author{}
\date{}

\begin{document}
\maketitle
\vspace{-2.2em}

Let $C$ be a finite-dimensional matrix with
\[
H:=\frac{C+C^*}{2}\geq 0,
\qquad
F(t):=\lVert e^{-Ct}\rVert_2^2.
\]
If $C$ is hypocoercive with index $m$, the known leading-order result gives
\begin{equation}\label{eq:known}
F(t)=1-c\,t^{2m+1}+O(t^{2m+2}),
\qquad c>0,
\end{equation}
as $t\downarrow0$.  The point of this note is that, when $m\geq1$, the
remainder in \eqref{eq:known} is automatically one order better: there is no
term of degree $2m+2$.  The first order at which a new coefficient can appear
is $2m+3$.

The mechanism is a parity property of the logarithmic norm.

\begin{lemma}[Odd logarithmic normal form]\label{lem:odd}
For every matrix $C$ there is a real-analytic function $\mu$, defined near
$0$, such that
\begin{equation}\label{eq:oddform}
\log F(t)=t\,\mu(t^2)
\end{equation}
for all sufficiently small $t>0$.  In particular, the right-hand Taylor
series of $\log F(t)$ contains only odd powers of $t$.
\end{lemma}

\begin{proof}
Set
\[
Q(t):=e^{-C^*t}e^{-Ct},
\qquad
Z(t):=\Log Q(t).
\]
Thus $Q(t)$ is positive definite, $Z(t)$ is Hermitian, and
$\log F(t)=\lambda_{\max}(Z(t))$.  The identity
\[
Q(-t)=e^{C^*t}Q(t)^{-1}e^{-C^*t}
\]
implies, by analytic functional calculus,
\begin{equation}\label{eq:similarity}
Z(-t)=-e^{C^*t}Z(t)e^{-C^*t}.
\end{equation}
Take the right polar decomposition $e^{C^*t}=U(t)R(t)$, where $U(t)$ is
unitary and $R(t)>0$.  Both sides of \eqref{eq:similarity} are Hermitian.
Consequently $R(t)Z(t)R(t)^{-1}$ is Hermitian, which gives
$R(t)^2Z(t)=Z(t)R(t)^2$.  Since $R(t)>0$, it follows that $R(t)$ commutes with
$Z(t)$, and hence
\begin{equation}\label{eq:unitary}
Z(-t)=-U(t)Z(t)U(t)^*.
\end{equation}

The polar unitary of an inverse is the adjoint, so $U(-t)=U(t)^*$.  Near
$t=0$, let $V(t):=U(t)^{-1/2}$ be the analytic principal inverse square root and put
$\widehat Z(t):=V(t)^*Z(t)V(t)$.  Since $V(-t)=U(t)V(t)$, equation
\eqref{eq:unitary} yields
\[
\widehat Z(-t)=-\widehat Z(t).
\]
Therefore $\widehat Z(t)=tB(t^2)$ for a real-analytic Hermitian matrix family
$B$.  By analytic Hermitian perturbation theory, the eigenvalues of $B(s)$
can be parametrized by real-analytic branches.  After shrinking the interval,
one branch $\mu(s)$ is maximal for every small $s>0$.  Since $\widehat Z(t)$
is unitarily similar to $Z(t)$,
\[
\log F(t)=\lambda_{\max}(Z(t))
=t\,\lambda_{\max}(B(t^2))=t\,\mu(t^2),
\]
which proves \eqref{eq:oddform}.
\end{proof}

\begin{proposition}[The missing even term]\label{prop:main}
Assume \eqref{eq:known} and $m\geq1$.  Then, for some real number $b$,
\begin{equation}\label{eq:refined}
F(t)=1-c\,t^{2m+1}+b\,t^{2m+3}
+O\!\left(t^{2m+5}+t^{4m+2}\right).
\end{equation}
In particular, the coefficient of $t^{2m+2}$ is zero and
\[
F(t)=1-c\,t^{2m+1}+O(t^{2m+3}).
\]
Thus $2m+3$ is the next \emph{admissible} order for an independent
correction.  The coefficient $b$ may itself vanish for a special matrix.
\end{proposition}

\begin{proof}
Since $\log(1+x)=x+O(x^2)$, equation \eqref{eq:known} shows that the leading
term of $\log F(t)$ is $-c\,t^{2m+1}$.  Lemma \ref{lem:odd} excludes every even
power, so for some $b\in\mathbb R$,
\[
\log F(t)=-c\,t^{2m+1}+b\,t^{2m+3}+O(t^{2m+5}).
\]
Exponentiating gives \eqref{eq:refined}: the first nonlinear contribution is
the square of the leading logarithmic term, hence has degree $4m+2$.  Because
$m\geq1$, this degree is strictly larger than $2m+2$, and the claimed
vanishing follows.
\end{proof}

For the first genuinely hypocoercive index, $m=1$, the statement is especially
transparent:
\[
F(t)=1-ct^3+bt^5+\frac{c^2}{2}t^6+O(t^7).
\]
Thus the fourth-order coefficient vanishes, the fifth-order coefficient is the
first potentially new one, and the sixth-order coefficient is already forced by $c$.

\paragraph{Why the coercive case is different.}
When $m=0$, one has $c=2\lambda_{\min}(H)>0$.  Lemma \ref{lem:odd} still gives
\[
\log F(t)=-ct+b t^3+O(t^5),
\]
but now exponentiation immediately produces an even term:
\begin{equation}\label{eq:coercive}
F(t)=1-ct+\frac{c^2}{2}t^2
+\left(b-\frac{c^3}{6}\right)t^3+O(t^4).
\end{equation}
Thus the assertion ``the remainder is one order higher'' is genuinely
hypocoercive: it relies not only on the oddness of $\log F$, but also on the
fact that its leading order is at least three.  In the coercive case the
quadratic term $c^2t^2/2=2\lambda_{\min}(H)^2t^2$ is unavoidable.

\vfill
\footnotesize
The leading law \eqref{eq:known} follows by squaring the finite-dimensional
short-time result of F.~Achleitner, A.~Arnold, and E.~A.~Carlen,
\emph{J. Differential Equations} \textbf{371} (2023), 83--115, Theorem~2.7.

\end{document}
